\documentclass[10pt,twocolumn]{article}

\usepackage[margin=0.9in]{geometry}
\usepackage{amsmath,amssymb,amsthm}
\usepackage{graphicx}
\usepackage{booktabs}
\usepackage{hyperref}
\usepackage{microtype}
\usepackage{cite}
\usepackage{enumitem}
\usepackage{float}
\usepackage{caption}
\usepackage{subcaption}
\usepackage{multirow}

\newtheorem{theorem}{Theorem}
\newtheorem{lemma}[theorem]{Lemma}
\newtheorem{corollary}[theorem]{Corollary}
\newtheorem{definition}{Definition}
\theoremstyle{remark}
\newtheorem{remark}{Remark}

\newcommand{\vrfi}{\textsc{VRFI}}
\newcommand{\vlr}{\textsc{VLR}}
\newcommand{\vadqn}{VA-DQN}
\newcommand{\vaobs}{VA-DQN-OBS}
\newcommand{\omegadqn}{$\omega$-DQN}
\newcommand{\omegavarn}{$\omega$-DQN-VarN}

% Agent naming:
% VA-DQN         = W10C  (2-head attention DQN, oracle burst)
% VA-DQN-OBS     = W12   (no burst time, observable features only)
% omega-DQN      = W14-omega (preference-conditioned, FiLM-style)
% omega-DQN-VarN = W15   (variable-N Poisson training)

\title{\textbf{Learning to Schedule: A Value-Aware Policy Discovery\\
Experiment with Provable Queue-Global Necessity}}

\author{Vedant Upadhyay\\
M.S.\ Computer Science, University of California Santa Cruz\\
Graduate Capstone Project, 2024--2025\\
\texttt{veupadhy@ucsc.edu}
\and
Scott Brandt\\
Department of Computer Science and Engineering\\
University of California Santa Cruz\\
\texttt{sbrandt@ucsc.edu}}

\date{March 2026}

\begin{document}

\maketitle

%% ============================================================
\begin{abstract}
%% ============================================================

CPU schedulers face a basic tension: you can optimize for throughput
or for fairness, but classical approaches struggle to do both at
once. Existing metrics like Jain's Fairness Index (JFI) and slowdown
variance (SDV) make this worse---they miss important types of
unfairness entirely.

We train a two-head attention Deep Q-Network (Value-Aware DQN,
\vadqn) on the Alibaba 2018 cluster trace. Each task gets an urgency
curve $V_i(d) = \mathrm{base}_i \cdot \max(\mathrm{floor}_i,\,
e^{-d/\tau_i})$, and we reward the agent for how much value it
preserves across the whole queue at each step.

We prove the Quantum Scoring Monotonicity Lemma: any scoring
function that looks at one task in isolation will always pick either
the smallest or the largest quantum---never something in between.
Three ablation experiments confirm this. Breaking out of this trap
requires looking at the whole queue, which is exactly what our
attention mechanism does.

\vadqn\ hits $21.71\pm0.64$\,s mean completion time across
3~independent training seeds (CV\,=\,2.9\%), beating Round Robin
with zero task starvation. Our best single checkpoint reaches
17.23\,s (N\,=\,1500 evaluation episodes).

We also introduce a new fairness metric, the Value-Rate Fairness
Index (\vrfi). Here is the key finding: ten tasks with identical
wait times score JFI\,=\,1.0 and SDV\,=\,0---both metrics say
``perfectly fair''---yet \vrfi\,=\,0.641 catches the real problem:
tasks with urgent value curves lose value twice as fast as slow ones.

A preference-conditioned variant (\omegadqn) beats the Multi-Level
Feedback Queue scheduler (MLFQ) on both throughput (20.60\,s vs.\
21.59\,s) and starvation (32.5\% vs.\ 36.0\%) at preference setting
$\omega_s=0.7$. No oracle information needed. No retraining to
switch modes.

\end{abstract}

%% ============================================================
\section{Introduction}
%% ============================================================

A CPU scheduler makes one decision, over and over: which waiting
process gets the CPU next, and for how long? The Linux Completely
Fair Scheduler (CFS) tracks how much CPU time each process has
received (via virtual runtime) and always schedules the one with
the least. This approximates fairness over time. However, it does
not account for job length, so it cannot optimize for metrics like
average completion time the way shortest-job-first scheduling can.
Short jobs still get CPU quickly due to frequent preemption, but
CFS does not preferentially complete them early the way SRPT would.

SRPT (Shortest Remaining Processing Time) solves the throughput
problem: always run the job closest to finishing. It minimizes mean
completion time in theory. The catch is that SRPT needs to know how
long each job will run, and in practice, schedulers never have that
information. It also starves long jobs.

We ask a different question. Can an agent learn a scheduling policy
from scratch---without copying an existing scheduler, without expert
demonstrations, without hand-tuned priority rules---just from
feedback about how long tasks have been waiting? We start by giving
the agent oracle burst time (an upper bound) and then strip it away
to see what the agent can do with only observable information.

This two-phase approach matters. It separates the question of what
the reward signal can teach from the question of what information is
actually available in production. If the agent degrades gracefully
when we remove oracle information, that tells us the reward design
is doing real work.

\subsection*{Agent Names}

We use descriptive names throughout.
\textbf{\vadqn} (formerly W10C): 2-head attention DQN with oracle
burst time.
\textbf{\vaobs} (formerly W12): same architecture, no burst time.
\textbf{\omegadqn} (formerly W14-$\omega$): preference-conditioned
with fixed affine modulation.
\textbf{\omegavarn} (formerly W15): trained on variable-length
Poisson arrival streams.

\subsection*{Contributions}

\begin{enumerate}[leftmargin=*, label=\arabic*., itemsep=2pt]

\item We document nine failure modes across four architectures---from
  tabular Q-learning to 2-head attention DQN---and diagnose each to
  a specific cause.

\item \vadqn\ achieves $21.71\pm0.64$\,s mean completion time
  (3~seeds, CV\,=\,2.9\%), with zero starvation. Best checkpoint:
  17.23\,s at N\,=\,1500 episodes. Two attention heads
  spontaneously specialize---one surveys the queue, one tracks the
  longest-waiting task---without any supervision.

\item We prove the Quantum Scoring Monotonicity Lemma and confirm it
  with three ablations: any value-drop scorer collapses to always
  picking $q_{\min}$ or always $q_{\max}$.

\item We find that default $\tau$ values (25\,s, 100\,s) are off by
  a factor of 44$\times$ and 8$\times$ from actual trace delays.
  Even with calibrated values, 93--100\% of ordering decisions are
  near-ties---a consequence of trace composition, not bad tuning.

\item \vrfi\ (Value-Rate Fairness Index) catches unfairness that JFI
  and SDV both miss. We show a concrete case: JFI\,=\,1.0,
  SDV\,=\,0, yet \vrfi\,=\,0.641.

\item A 2$\times$2 ablation shows that reward structure matters more
  than state observability. Swapping out the value-delta reward
  collapses SRPT agreement to 11.3\%. Removing oracle burst costs
  roughly 4.4\,s relative to the best checkpoint.

\item \omegadqn\ improves both throughput and starvation relative to
  MLFQ at $\omega_s=0.7$. One model, adjustable at runtime, no
  retraining.

\item \omegavarn\ confirms the MCT advantage holds on variable-length
  queues (mean $-1.40$\,s vs.\ MLFQ, 3~seeds).

\end{enumerate}

%% ============================================================
\section{Background}
%% ============================================================

\subsection{Why CFS Falls Short}

CFS schedules the process with the smallest virtual runtime,
approximating proportional CPU fairness across tasks. Fairness does
not optimize for latency though: unlike shortest-job-first, CFS does
not prioritize short tasks for early completion. Short jobs still get
CPU quickly due to frequent preemption, but CFS does not explicitly
favor them.

Since Linux 6.6, the default scheduler uses an EEVDF-based
model~\cite{eevdf}, which adds virtual deadlines and eligibility
times to more precisely control when tasks become runnable under fair
scheduling. Our approach differs from both: instead of virtual
runtime or deadlines, we model per-task urgency curves directly. That
shifts the goal from time-fairness to value-fairness---tasks that
lose value faster get scheduled sooner.

\subsection{What Prior RL Schedulers Do}

KernelOracle~\cite{kernelOracle} trains an LSTM to predict which
task CFS will select next---it learns to mimic an existing policy,
not discover a new one. ALPS~\cite{alps} approximates SRPT by
simulating optimal schedules on historical workloads and enforcing
the resulting priorities via eBPF, reporting a 57.2\% latency
reduction in serverless environments. RL-based schedulers like
RLScheduler~\cite{rlscheduler}, Double DQN approaches~\cite{doubleDQN},
and SmartOS~\cite{smartOS} learn policies using workload statistics,
job classes, or multi-resource signals.

Most prior approaches rely on historical runtime information,
workload-specific priors, or reference policies. Our approach uses
only simple, online-computable features---remaining burst, wait time,
and arrival indicators---without offline profiling or job-class
metadata. One caveat: remaining burst time is oracle information in
our experiments, unavailable in real kernels without estimation.
Section~\ref{sec:2x2} measures the exact cost of removing it:
roughly 4.4\,s MCT on the Alibaba 2018 trace.

\subsection{Attention for Queue-Global Decisions}

Standard DQN scores each action independently. That breaks down for
scheduling because the right choice for task $i$ depends on what
every other task in the queue needs. We use scaled dot-product
attention~\cite{vaswani2017} so each task sees the full queue before
getting scored. Section~\ref{sec:theory} explains formally why this
matters.

\subsection{Production Systems Hit the Same Wall}

A 2026 public preview of cloud batch scheduler fairness added a
server-side backlog tracker because per-task metadata alone was not
enough. That is exactly what our Corollary~\ref{cor:global} proves:
any scoring function that only looks at one task at a time will be
degenerate. Queue-global state is not a nice-to-have---it is
necessary.

%% ============================================================
\section{Theoretical Framework}
\label{sec:theory}
%% ============================================================

\subsection{Value Curves}

We assign each task $i$ a value that decays as it waits:
\begin{equation}
  V_i(d) = \mathrm{base}_i \cdot
  \max\!\bigl(\mathrm{floor}_i,\; e^{-d/\tau_i}\bigr)
  \label{eq:vcurve}
\end{equation}
where $d$ is how long the task has waited, $\tau_i$ controls how
fast it loses value, $\mathrm{base}_i$ is its starting value, and
$\mathrm{floor}_i$ is the minimum it can reach.

We use two profiles: \textbf{steep} ($\tau\approx25$\,s,
floor$\,=0.2$, for interactive/IO tasks) and \textbf{smooth}
($\tau\approx100$\,s, floor$\,=0$, for batch tasks). These
assignments are human choices---they encode our prior about task
urgency and are not learned from the trace.

At each scheduling step, we pick the task losing value fastest right
now:
\begin{equation}
  i^* = \arg\max_i \left(-\frac{dV_i}{dt}\right)
  = \arg\max_i \frac{\mathrm{base}_i}{\tau_i}\,
  e^{-d_i/\tau_i}
  \label{eq:policy}
\end{equation}
This is greedy minimization of instantaneous value loss---the same
spirit as SRPT, but generalized to arbitrary urgency curves rather
than just job length.

The floor prevents tasks from becoming irrelevant. Combined with
diminishing gradients---a long-waiting task eventually hits its
floor and stops competing, freeing the scheduler to serve
others---the design avoids indefinite postponement in practice.
This is not a hard fairness guarantee; Section~\ref{sec:pbrs}
discusses the limits of soft starvation prevention in detail.

In practice, the RL agent does not compute
equation~\eqref{eq:policy} directly. It learns to approximate this
selection rule from the value-delta reward signal, which encodes the
same gradient information implicitly at each training step. The
benefit of RL over the analytic rule is generalization: the analytic
rule requires knowing all task parameters perfectly, while the agent
learns from observed features and adapts to real trace distributions.

\subsection{The Potential Surface}

We model the queue as a potential surface:
\begin{equation}
  \Phi(S) = \sum_i V_i(d_i)
  \label{eq:potential}
\end{equation}
where each task contributes value that decays with waiting time. The
scheduler greedily minimizes instantaneous loss in this total
value---it does not solve a long-horizon objective, just picks the
best action right now.

For each task, the rate of value decay is:
\begin{equation}
  \frac{dV_i}{dt} = -\frac{\mathrm{base}_i}{\tau_i}\,
  e^{-d_i/\tau_i}
  \quad\text{(non-plateau)};\qquad
  = 0\quad\text{(plateau)}
  \label{eq:decay}
\end{equation}

At each decision point, we select the task with the largest
magnitude of value decay:
\begin{equation}
  i^* = \arg\max_i \left|\frac{dV_i}{dt}\right|
  = \arg\max_i \frac{\mathrm{base}_i}{\tau_i}\,
  e^{-d_i/\tau_i}
  \label{eq:selection}
\end{equation}
This picks the task whose delay incurs the highest immediate cost to
the system.

The decision is inherently comparative: priority depends on relative
decay rates across all tasks, not on any task's absolute value in
isolation. A task with a high remaining value but slow decay loses to
a task with lower remaining value but fast decay. This is why the
selection rule requires observing the whole queue---you need
everyone's decay rate to know whose is largest.

Tasks on the plateau have zero decay and drop out of contention until
competing tasks' own decay rates diminish. This is not permanent
starvation---once higher-decay tasks are served and their gradients
shrink, plateau tasks become competitive again.

\subsection{Degeneracy in Local Quantum Scoring}
\label{sec:lemma}

Once you pick a task, you also pick a quantum---how long to run it.
Two natural ways to score quanta for task $i$ are:
\begin{align}
  f_{\mathrm{rate}}(q) &= \frac{V_i(d_i)-V_i(d_i+q)}{q}
  \quad\text{(value preserved per second)}
  \label{eq:frate}\\
  f_{\mathrm{total}}(q) &= V_i(d_i)-V_i(d_i+q)
  \quad\text{(total value preserved)}
  \label{eq:ftotal}
\end{align}
Both look only at task $i$'s own curve. We show this is a
fundamental limitation.

\begin{lemma}[Quantum Scoring Monotonicity]
\label{lem:mono}
For any $\tau>0$, $\mathrm{base}>0$,
$\mathrm{floor}\in[0,1)$, $d\geq0$:
\begin{enumerate}[label=(\alph*)]
  \item $f_{\mathrm{rate}}(q)$ strictly decreases as $q$ grows.
  \item $f_{\mathrm{total}}(q)$ strictly increases as $q$ grows
    (when $d+q$ stays below the plateau).
\end{enumerate}
\end{lemma}

\begin{proof}
\textbf{(a):} Write $f_{\mathrm{rate}}(q)=
\mathrm{base}\cdot e^{-d/\tau}\cdot h(q)$ where
$h(q)=(1-e^{-q/\tau})/q$. Differentiating gives
$h'(q)=[(q/\tau+1)e^{-q/\tau}-1]/q^2$. Set
$u=q/\tau>0$. The numerator $(u+1)e^{-u}-1$ is
negative for all $u>0$ because $e^u>1+u$ by strict
convexity. So $h'(q)<0$ everywhere.\hfill$\square$

\textbf{(b):} Below the plateau,
$f_{\mathrm{total}}(q)=\mathrm{base}\cdot e^{-d/\tau}
\cdot(1-e^{-q/\tau})$. Since $1-e^{-q/\tau}$ grows
with $q$, so does $f_{\mathrm{total}}$.\hfill$\square$

\textbf{At the plateau:} Once a task hits its floor,
$f_{\mathrm{total}}(q)=0$ for every $q$. The task
contributes no marginal value signal, though it
remains in the system and can still be scheduled.
\end{proof}

\begin{corollary}[Degeneracy of Local Value Scorers]
\label{cor:global}
For the two natural value-preservation objectives above, any policy
based solely on task-local information selects only extreme quanta:
$f_{\mathrm{rate}}$ always favors $q_{\min}$; $f_{\mathrm{total}}$
always favors $q_{\max}$. Neither can produce balanced quantum
selection. Other local rules not based on value decay may behave
differently---this result applies specifically to
$f_{\mathrm{rate}}$ and $f_{\mathrm{total}}$.
\end{corollary}

This limitation runs deeper than it first appears. Quantum selection
is not really a single-task decision---it is a multi-task tradeoff.
Running task $i$ for a long quantum means every other waiting task
loses value during that time. The right quantum for task $i$
therefore depends on what the rest of the queue needs right now,
information that neither $f_{\mathrm{rate}}$ nor $f_{\mathrm{total}}$
can see.

Local scoring is insufficient because it ignores these cross-task
interactions. Our model addresses this by letting each task condition
on the full queue before any scoring happens---each task attends to
every other task, so the quantum decision reflects the cost to the
whole queue, not just the benefit to one task.

%% ============================================================
\section{Method}
\label{sec:method}
%% ============================================================

\subsection{Environment}

\textbf{Synthetic (architecture search, Weeks 1--9).}
We use small synthetic instances (5 processes per episode) to enable
rapid architecture search while preserving non-trivial scheduling
interactions. Arrivals are sampled from $\{0,2,5,8,10\}$\,ms and
burst lengths from $U[1,60]$\,ms, resampled each episode. This
simple uniform distribution does not reflect real workload skew---it
provides a controlled setting for model selection before moving to
real traces. The action space is 15 choices (5 processes $\times$
3 quanta: 1\,ms, 5\,ms, 20\,ms), spanning short, medium, and long
scheduling intervals. Invalid actions are masked at every step.

\textbf{Real-trace (Weeks 8--15).} We use the Alibaba 2018 cluster
trace (\texttt{batch\_task} table, 14.3M tasks). We keep only tasks
with duration up to 47.0\,s (the 75th percentile) to reduce extreme
long-tail effects during training and maintain tractable
within-episode burst ratios---without this filter, the p95/p50 burst
ratio reaches 39.7$\times$, which destabilizes early training. This
filtering focuses evaluation on the majority workload regime;
behavior on the upper quartile is left as future work. The filtered
set splits randomly into 8.18M training and 2.04M test tasks. Quanta
for real-trace agents are 0.5\,s, 2.0\,s, and 8.0\,s, spanning
short, medium, and long scheduling intervals. For larger queues, the
policy operates over a variable number of tasks using
attention-based representations (see Section~\ref{sec:arch}).

\subsection{Reward}

\textbf{MCT reward (Weeks 1--10).} At each step:
\begin{equation}
  R = -n_{\mathrm{active}} \times q_{\mathrm{actual}}
\end{equation}
where $n_{\mathrm{active}}$ counts arrived but unfinished processes
and $q_{\mathrm{actual}}=\min(q,\,\mathrm{remaining})$. Summing
this over an episode measures the area under the active-job curve,
which equals total flow time across all jobs---minimizing it
therefore minimizes average turnaround time~\cite{littlesLaw}. We
use no bonus terms or shaping.

\textbf{Value-delta reward (real-trace variants).}
\begin{equation}
  R = \tfrac{1}{20}\sum_{i\in\mathrm{runnable}}
  \bigl[V_i(d_i+q) - V_i(d_i)\bigr]
\end{equation}
Because task value decays with waiting, this reward is always
negative---it measures the value lost across the whole queue during
the chosen interval. Tasks with steep current decay incur larger
losses, so the agent learns to take actions that reduce those losses.
This penalizes value loss and induces urgency-sensitive behavior as a
consequence.

\textbf{\omegadqn\ reward.} We combine value preservation with a
starvation penalty:
\begin{equation}
  R = \omega_{\mathrm{vd}} \cdot R_{\mathrm{vd}}
  + \omega_s \cdot R_{\mathrm{ss}}
\end{equation}
where the starvation signal is:
\begin{equation}
  R_{\mathrm{ss}} = -\frac{1}{N}\sum_i
  \max(0,\, t_{\mathrm{starve},i} - 50)\,/\,50
\end{equation}
This signal is dense and activates at every step. In preliminary
experiments, binary starvation penalties were less stable---the
proportional signal gives the agent a gradient to follow rather than
a cliff to avoid.

\subsection{What Goes Into the State}

\textbf{Oracle phase (\vadqn).} Feature \texttt{off+0} is true
remaining burst time---information unavailable to real schedulers.
We include it deliberately to estimate an upper bound on achievable
performance under our model and reward formulation, so we can
quantify the cost of removing it.

\textbf{Observable phase (\vaobs, \omegadqn, \omegavarn).} We
replace oracle burst with seven features a practical scheduler can
observe at runtime:

\begin{itemize}[itemsep=2pt]
  \item \textbf{time in queue} --- total elapsed time since arrival
  \item \textbf{cumulative wait time} --- time spent waiting,
    excluding any CPU time already received
  \item \textbf{time since last execution} --- starvation signal;
    grows when a task is repeatedly skipped
  \item \textbf{value loss rate} --- derived from our predefined
    value curves, not directly observable in a real kernel
  \item \textbf{CPU request} --- planned CPU usage from the trace
  \item \textbf{memory request} --- planned memory usage from the
    trace
  \item \textbf{arrival flag} --- whether the task has arrived yet
\end{itemize}

This gives 35 dimensions (5 processes $\times$ 7 features). For
larger queues, the same per-task features feed into
attention-based representations that scale to variable queue sizes
(Section~\ref{sec:arch}).

One feature warrants a caveat: value loss rate is computed from our
designed value curves and requires knowing $\tau_i$ and
$\mathrm{base}_i$ per task. It is not a raw kernel observable. We
include it because it encodes urgency compactly, but we acknowledge
it depends on the value model assumptions stated in
Section~\ref{sec:theory}.

The ablation in Section~\ref{sec:2x2} measures the cost of this
entire swap: removing oracle burst increases MCT by about 4.4\,s on
the 2018 trace, roughly a 20\% increase over the oracle-equipped
baseline.

\subsection{Architecture}
\label{sec:arch}

\begin{figure}[t]
  \centering
  \includegraphics[width=\columnwidth]{fig3_architecture.pdf}
  \caption{\omegadqn\ architecture. The preference scalar $\omega_s$
  enters at two points: it modulates attention sharpness before the
  attention operation
  ($Q_{\mathrm{cond}}=Q\times(1+\omega_s)$) and scales the context
  vector after attention
  ($c_{\mathrm{cond}}=c\times(1+\omega_s)+(1-\omega_s)$). The MLP
  scores each candidate from a 24-dim input: task features (7) +
  context (16) + $\omega_s$ (1).}
  \label{fig:arch}
\end{figure}

\textbf{\vadqn.} Two attention heads ($D_{\mathrm{head}}=8$) over
the 5-process encodings, using scaled dot-product
attention~\cite{vaswani2017}. We project the concatenated head
outputs to a 16-dim context vector, then score each candidate with
an MLP: $[\mathbf{e}_i(7)\,\|\,q_t(1)]\to64\to32\to1$, where $q_t$
is the normalized quantum tier. The quantum enters the MLP directly,
so the agent learns separate Q-values for each (task, quantum)
pair---for example, (task~2, $q=0.5$\,s) and (task~2, $q=8.0$\,s)
receive independent scores. Total: 4,145 parameters. Training uses a
Double DQN objective~\cite{mnih2015}: $\varepsilon$-greedy decay
$1.0\to0.05$, Adam ($\mathrm{lr}=10^{-4}$), $\gamma=0.99$, replay
buffer 10,000, batch 32, target sync every 5 episodes.

\omegavarn\ differs in one important way: the quantum tier does not
enter the MLP. Instead, one task-level score is tiled across all
three quanta, reducing quantum selection to a random choice among
valid options for the chosen task. This isolates the contribution of
explicit quantum modeling: \vadqn\ learns distinct Q-values per
(task, quantum) pair, while \omegavarn\ removes this capacity
entirely. We report this difference explicitly for reproducibility.

\textbf{Fixed affine conditioning (\omegadqn).} Let
$\omega_{\mathrm{mct}} = 1 - \omega_s$, so the two scalars sum to 1.
The preference scalar $\omega_s\in[0,1]$ modulates two points in the
forward pass:
\begin{align}
  Q_{\mathrm{cond}} &= Q\cdot(1+\omega_s)
  \label{eq:film1}\\
  c_{\mathrm{cond}} &= c\cdot(1+\omega_s) + \omega_{\mathrm{mct}}
  \label{eq:film2}
\end{align}
The first operation modulates attention sharpness and
sensitivity---higher $\omega_s$ amplifies the query magnitudes,
producing more peaked attention distributions. The second shifts the
context vector toward the starvation objective as $\omega_s$ grows.
We use fixed scaling rather than learned parameters, so $\omega_s$
acts as a direct control knob the network cannot learn to ignore.
This differs from standard FiLM~\cite{film2018} in that the affine
coefficients are not learned functions of the conditioning
input---the design choice prioritizes interpretability and control
over flexibility.

\textbf{Preference sampling.} For the first 8,000 episodes we draw
$\omega_s \sim U(0,1)$. After that we switch to a
Beta(0.5,\,0.5)-like distribution, approximated as
$\omega_s=\sin^2(u\pi/2)$, $u\sim U[0,1]$. This concentrates
samples near 0 and 1, emphasizing the extreme policies before
interpolation and improving stability in preference-conditioned
training. We store $\omega_s$ alongside each transition in the
replay buffer and never resample it at update time---resampling would
make the Q-targets inconsistent with the transitions that generated
them.

\subsection{Baselines}

SRPT oracle (theoretical ceiling, needs true burst lengths), Round
Robin (no ordering, fixed quantum), MLFQ (3 queues, quanta
$\{4,8,16\}$\,s, aging threshold 50\,s), CFS-lite (vruntime, 1\,s
granularity, 10\,s target latency), EEVDF-lite (CFS-lite with
per-task eligible times).

%% ============================================================
\section{Results}
\label{sec:results}
%% ============================================================

\subsection{\vadqn: Main Result}

\begin{figure}[t]
  \centering
  \includegraphics[width=\columnwidth]{fig2_learning_curve.pdf}
  \caption{\vadqn\ learning curve. The gray band shows replicated
  training across 3 seeds: $21.71\pm0.64$\,s (CV\,=\,3.0\%). The
  best checkpoint hits 17.23\,s (N\,=\,1500 test episodes). The two
  variance numbers measure different things: $\pm0.64$\,s is
  seed-level variance across training runs; $\pm11.38$\,s is
  episode-level variance reflecting workload diversity, not policy
  instability.}
  \label{fig:learning}
\end{figure}

Table~\ref{tab:main} shows the full progression. Across
3~independent seeds, \vadqn\ achieves $\mathbf{21.71\pm0.64}$\,s
(CV\,=\,2.9\%). Mean performance is comparable to Round Robin
(21.31\,s), but the learned policy is substantially more
consistent---RR's episode-level std is $\pm14.03$\,s vs.\
$\pm0.64$\,s seed-level variance for \vadqn. The best single
checkpoint reaches 17.23\,s on N\,=\,1500 test episodes; we report
the replicated mean as the primary result to avoid selection bias
over stochastic episodes (Figure~\ref{fig:learning}). The 4.5\,s gap
between the best checkpoint and the replicated mean reflects this
selection effect rather than training instability---seed-level
variance stays low at 3.0\%. SRPT oracle achieves 15.20\,s,
indicating roughly 6.5\,s of remaining headroom.

\begin{table}[h]
\centering
\caption{Agent results on Alibaba 2018 test split.
$^\dagger$Best checkpoint, N\,=\,1500 eval.
Replicated training: $21.71\pm0.64$\,s (3~seeds, CV\,=\,3.0\%).}
\label{tab:main}
\small
\begin{tabular}{lrr}
\toprule
Agent & MCT (s) & vs RR \\
\midrule
SRPT oracle & 15.20 & $-$6.11 \\
\vadqn\ (best ckpt$^\dagger$) & $17.23\pm11.38$ & $-$4.08 \\
W9 (1-head) & $19.22\pm11.58$ & $-$2.09 \\
\textbf{\vadqn\ (replicated)} & $\mathbf{21.71\pm0.64}$ & $\sim$0 \\
Round Robin & $21.31\pm14.03$ & baseline \\
\bottomrule
\end{tabular}
\end{table}

We observe consistent qualitative patterns in attention behavior
across training runs: one head tends to attend broadly across the
queue, while the other often concentrates on long-waiting tasks. We
do not claim a mechanistic explanation---these are qualitative
observations from inspecting attention weights, not a controlled
analysis. Agreement with SRPT in task selection reaches 71.9\%
overall, suggesting the agent approximates urgency-based ordering
without being explicitly told to do so.

\subsection{Ablation: The Lemma Holds Empirically}
\label{sec:theory-confirms}

\begin{table}[h]
\centering
\caption{Quantum degeneracy ablation confirming Lemma~\ref{lem:mono}.
Total Value is the cumulative sum of $V_i(d_i)$ across all tasks and
timesteps in an episode---higher means more value preserved over
time. Local scorers are evaluated as direct implementations of their
objectives without additional tuning.}
\label{tab:abl}
\small
\begin{tabular}{lrrr}
\toprule
Variant & $q_{\min}$\% & $q_{\max}$\% & Total Value \\
\midrule
$f_{\mathrm{rate}}$ scorer & 100 & 0 & 250.55 \\
$f_{\mathrm{total}}$ scorer & 0 & 100 & 159.21 \\
Ordering only ($q_{\mathrm{med}}$ fixed) & --- & --- & 15.37 \\
\textbf{\vadqn\ (attention)} & mixed & mixed & 241.92 \\
\bottomrule
\end{tabular}
\end{table}

Table~\ref{tab:abl} confirms the lemma directly. Both local scorers
collapse to a single extreme---exactly as predicted. $f_{\mathrm{rate}}$
picks $q_{\min}$ 100\% of the time; $f_{\mathrm{total}}$ picks
$q_{\max}$ 100\% of the time. Neither deviates once across all
evaluated episodes.

\vadqn\ is the only policy among those evaluated that consistently
picks mixed quanta. Its total value (241.92) is lower than
$f_{\mathrm{rate}}$ (250.55), but $f_{\mathrm{rate}}$ achieves this
by always selecting the smallest quantum regardless of queue
state---a degenerate strategy that happens to score well on this
metric while being useless in practice. \vadqn\ achieves comparable
total value without collapsing to an extreme.

The ``ordering only'' variant is the most revealing result. Even
with correct task ordering, fixing a middle quantum leads to
near-complete smooth-task starvation, as measured by sustained low
value accumulation for smooth tasks across the episode (total value
15.37 vs.\ 241.92). This shows that task selection and quantum
selection are not independent problems---getting the order right is
not enough. Effective scheduling requires joint reasoning over both
decisions simultaneously, which is what queue-global attention
enables.

\subsection{Why $\tau$ Calibration Matters---and Why It Is Not Enough}

Default $\tau$ values are 25\,s (steep) and 100\,s (smooth), while
actual median delays in the trace are 5,279\,s and 38,665\,s
respectively. At these scales, $e^{-5279/25}$ is on the order of
$10^{-92}$---effectively zero for scheduling purposes. The value
gradient is numerically negligible before any task ever runs, so the
scheduler has no signal to act on regardless of which task it picks.

Increasing $\tau$ to 900\,s (steep) and 800\,s (smooth) restores
non-negligible gradients. But it does not resolve the decision
quality problem. The gradient ratio between two tasks is:
\begin{equation}
  \frac{\nabla\Phi_i}{\nabla\Phi_j} =
  e^{-(d_i-d_j)/\tau}
\end{equation}
When $\tau$ is large relative to the delay difference $|d_i - d_j|$,
this ratio approaches 1---the two tasks are nearly indistinguishable
by gradient alone. In our trace, 92\% of tasks are smooth with
similar $\tau$ values, so most queues consist of tasks whose
gradients are nearly identical.

We define a near-tie as any pair of candidates whose gradient ratio
falls within 1\% of each other
($|\nabla\Phi_i - \nabla\Phi_j| / \nabla\Phi_i < 0.01$). Under this
definition, 93--100\% of ordering decisions are near-ties across all
tested policies, regardless of $\tau$ calibration. This suggests the
issue is primarily driven by trace composition---a workload dominated
by similar smooth tasks---rather than miscalibration of $\tau$ alone.
Better calibration is necessary but not sufficient.

\subsection{VRFI: A Fairness Metric Sensitive to Urgency Differences}
\label{sec:vrfi}

\begin{figure}[t]
  \centering
  \includegraphics[width=\columnwidth]{fig4_vrfi_falsification.pdf}
  \caption{The falsification case for JFI and SDV. Ten tasks all wait
  the same amount ($d=50$\,s). Both classical metrics report perfect
  fairness. But steep tasks (orange) lose value at approximately
  twice the rate of smooth tasks (blue). \vrfi\,=\,0.641 detects
  this disparity; JFI and SDV do not.}
  \label{fig:vrfi}
\end{figure}

\begin{definition}[\vlr\ and \vrfi]
For task $i$ with delay $d_i$, define the value loss rate as the
average loss over the observed delay:
\begin{equation}
  \vlr_i = \frac{V_i(0)-V_i(d_i)}{d_i}
\end{equation}
We use average loss rate rather than instantaneous decay for
stability under discrete scheduling intervals; for small $d_i$ this
approximates the initial decay rate $\mathrm{base}_i/\tau_i$.

The Value-Rate Fairness Index is then:
\begin{equation}
  \vrfi = 1-\frac{\sigma(\vlr)}{\mu(\vlr)} =
  1-\mathrm{CV}(\vlr)
\end{equation}
\vrfi\,=\,1 means every task loses value at the same rate. As
inequality across tasks increases, \vrfi\ decreases without bound.
\end{definition}

\textbf{The falsification case.} Ten tasks, all waiting $d=50$\,s:
five steep ($\tau=25$\,s) and five smooth ($\tau=100$\,s). JFI sees
equal turnarounds and reports 1.0. SDV reports zero variance under
identical delays. But steep tasks lose value at approximately
$2\times$ the rate of smooth tasks---\vrfi\,=\,0.641 detects this
disparity while both existing metrics miss it entirely
(Figure~\ref{fig:vrfi}).

Table~\ref{tab:fair} compares five policies. \vrfi\ is defined only
for policies operating under a value-based framework where task
urgency curves are specified; it is not applicable to RR, CFS-lite,
or FCFS, which carry no value model. Among the policies where it
applies, the value-aware policy achieves the highest \vrfi\ (0.712)
alongside zero starvation. Round Robin wins JFI precisely because it
treats all tasks identically---but uniform treatment ignores urgency
differences, which is the failure mode \vrfi\ is designed to expose.

\begin{table}[h]
\centering
\caption{Five-policy fairness comparison (300-task slice). \vrfi\
applies only to value-aware policies where urgency curves are
defined.}
\label{tab:fair}
\small
\begin{tabular}{lrrrr}
\toprule
Policy & JFI & SDV & Starv. & \vrfi \\
\midrule
Value-aware & 0.934 & 566 & \textbf{0} & 0.712 \\
MLFQ & 0.943 & 2,193 & 0 & 0.420 \\
CFS-lite & 0.916 & 6,138 & 1 & n/a \\
Round Robin & \textbf{0.953} & 15,819 & 5 & n/a \\
FCFS & 0.774 & 12.7M & 27 & n/a \\
\bottomrule
\end{tabular}
\end{table}

\subsection{What We Learn From the 2$\times$2 Ablation}
\label{sec:2x2}

\begin{table*}[t]
\centering
\caption{2$\times$2 ablation (Alibaba 2018, mean$\pm$std, 3~seeds).
$^*$\vadqn\ best checkpoint N\,=\,1500; replicated mean is 21.71\,s.
All other agents report replicated means.}
\label{tab:2x2}
\small
\begin{tabular}{llrr}
\toprule
What we changed & Agent & MCT (s) & SRPT\,(\%) \\
\midrule
Nothing (baseline)   & \vadqn\ (ckpt$^*$) & $17.23\pm11.38$ & $45.7\pm3.0$ \\
Nothing (replicated) & \vadqn             & $21.71\pm0.64$  & $45.7\pm3.0$ \\
Removed burst        & \vaobs             & $21.67\pm0.42$  & 56.6 \\
Changed reward       & W11d               & $26.35\pm1.93$  & 40.6 \\
Changed both         & W11/W11b           & $21.05\pm0.72$  & 54.7 \\
\bottomrule
\end{tabular}
\end{table*}

We ran four agents that each change one or both factors
(Table~\ref{tab:2x2}). All comparisons below use replicated means
unless stated otherwise.

\textbf{Removing burst time costs roughly 0.04\,s on the replicated
mean.} \vaobs\ removes oracle burst and replaces it with observable
features, moving from 21.71\,s to 21.67\,s---effectively the same.
Relative to the best checkpoint (17.23\,s), the apparent gap is
4.4\,s, but this reflects checkpoint selection bias rather than a
clean causal effect. The honest reading is that removing oracle burst
has modest impact when the reward remains value-delta.

\textbf{Changing the reward has a larger effect.} W11d keeps oracle
burst but replaces value-delta with a composite reward. MCT rises to
26.35\,s and SRPT agreement drops to 40.6\%. In this setting, reward
structure has a larger impact than state observability.

\textbf{Without burst, reward choice has relatively small effect.}
\vaobs\ (21.67\,s) and W11/W11b (21.05\,s) differ by only 0.6\,s.
Once oracle burst is removed, changing the reward has limited
additional effect, suggesting the state representation becomes the
dominant constraint at this performance level.

\textbf{Oracle burst combined with the wrong reward degrades
performance.} W11d (burst kept, reward changed) is worse than
W11/W11b (both changed): 26.35\,s vs.\ 21.05\,s. When combined with
a mismatched reward, precise state information may amplify
optimization of the wrong objective rather than help.

One result worth noting: \vaobs\ achieves higher SRPT agreement
(56.6\%) than \vadqn\ (45.7\%) yet has worse MCT. Higher agreement
with SRPT task ordering does not directly translate to lower
completion time---quantum selection and reward alignment also play a
role.

Why does value-delta work? It couples efficiency and urgency into a
single signal. A task early in its decay curve contributes a large
negative delta---the agent learns to finish it quickly as a
consequence of maximizing reward. When that signal is split into
separate additive components, the coupling breaks and performance
falls. The W11 series did not prove this mechanism directly, but the
pattern is consistent across all four conditions.

\subsection{Why Starvation Is Hard to Fix With Reward Shaping}
\label{sec:pbrs}

We evaluate potential-based reward shaping (PBRS) using a starvation
potential:
\begin{equation}
  \varphi(s) = -\lambda \cdot
  \max_i(t_{\mathrm{starve},i})
\end{equation}
applied as $r' = r + \gamma\varphi(s') - \varphi(s)$, consistent
with the PBRS framework. We test five $\lambda$ values spanning weak
to strong penalties (0.01--0.50). Starvation stays between 46\% and
57\% in every case. No configuration crosses the 36\% MLFQ threshold
while keeping MCT below 21.59\,s simultaneously.

MLFQ addresses starvation through an aging rule: any task waiting
beyond 50\,s is promoted to the highest priority queue, providing a
hard scheduling guarantee. Reward shaping provides only a soft
incentive. In our experiments, the starvation penalty is dominated by
the value-delta signal, which produces larger reward magnitudes per
step---the agent consistently prioritizes throughput across all
tested $\lambda$ values. Continuous shaping was insufficient to
replicate the behavior of a hard aging rule in this setting.

Two points worth noting. First, we only tested $\lambda$ up to 0.50;
larger values might reduce starvation further but would likely
destabilize MCT. Second, a different shaping design---one that more
directly penalizes the scheduling decision rather than the
state---might behave differently. We leave that direction open.

\subsection{The Preference-Conditioned Agent (\omegadqn)}

\begin{figure}[t]
  \centering
  \includegraphics[width=\columnwidth]{fig1_pareto_frontier.pdf}
  \caption{\omegadqn\ Pareto frontier. The shaded region shows
  $\omega_s \in [0.6, 0.8]$, where the policy dominates MLFQ on
  both MCT and starvation simultaneously. We recommend $\omega_s=0.7$
  (gold star): MCT\,=\,20.60\,s, starvation\,=\,32.5\%. The
  frontier is non-monotonic: pushing $\omega_s$ to 1.0 produces
  worse starvation than 0.7.}
  \label{fig:pareto}
\end{figure}

Instead of fixing a single tradeoff at training time, we train one
policy that accepts a preference scalar $\omega_s\in[0,1]$ at
runtime. Lower values bias the policy toward minimizing completion
time; higher values bias it toward reducing starvation. Intermediate
values interpolate between both objectives. No retraining needed to
switch behavior.

The ratio of gradient norms from the two reward components stayed
between 0.92 and 1.12 throughout 30,000 training episodes,
indicating that neither objective dominated the other during
training. Unlike the earlier composite reward experiments in
Section~\ref{sec:pbrs}, we observe no collapse in policy behavior
across the full training run.

Table~\ref{tab:omega} shows the full $\omega_s$ sweep.
$\omega_s\in[0.6, 0.8]$ dominates MLFQ on both metrics
simultaneously. Results are single-seed; seed-level variance for
this agent is similar to that reported for \vadqn\ in
Table~\ref{tab:main}.

\begin{table}[h]
\centering
\caption{\omegadqn\ Pareto frontier.
$\omega_s\in[0.6,0.8]$ dominates MLFQ on both MCT and starvation.}
\label{tab:omega}
\small
\begin{tabular}{rrrr}
\toprule
$\omega_s$ & MCT (s) & Starv.\,(\%) & vs MLFQ \\
\midrule
0.0 & 20.51 & 44.0 & MCT\,$\checkmark$, Starv.\,$\times$ \\
0.5 & 20.73 & 38.0 & MCT\,$\checkmark$, Starv.\,$\times$ \\
\textbf{0.6} & \textbf{20.49} & \textbf{35.0} & \textbf{Both\,$\checkmark$} \\
\textbf{0.7} & \textbf{20.60} & \textbf{32.5} & \textbf{Both\,$\checkmark$} \\
\textbf{0.8} & \textbf{20.43} & \textbf{34.5} & \textbf{Both\,$\checkmark$} \\
0.9 & 20.45 & 38.5 & MCT\,$\checkmark$, Starv.\,$\times$ \\
1.0 & 20.36 & 38.5 & MCT\,$\checkmark$, Starv.\,$\times$ \\
\midrule
MLFQ & 21.59 & 36.0 & baseline \\
\bottomrule
\end{tabular}
\end{table}

The non-monotonic shape of the frontier is notable. Increasing
$\omega_s$ beyond 0.7 leads to worse starvation outcomes rather than
better ones. This suggests that some throughput pressure is necessary
to avoid excessive delay accumulation across tasks---when the policy
deprioritizes throughput entirely, secondary queueing effects appear
to worsen starvation rather than eliminate it. This observation is
consistent with classical schedulers using non-zero aging thresholds
rather than immediate promotion.

$\omega_s$ acts as a runtime control knob, analogous to system-level
tuning parameters like \texttt{nice} values. The difference is that
those knobs are hand-engineered heuristics. \omegadqn\ learns the
same control surface from data.

\subsection{Does It Generalize to Variable Queue Depths? (\omegavarn)}

\vadqn\ was trained on exactly five processes. When we test it on a
Poisson arrival stream where queue depth varies---truncating to five
slots when the queue grows longer---it matches MLFQ at low load
($\rho=0.3$) but loses at moderate load ($\rho=0.7$). The bottleneck
is the five-slot window: once more than five tasks are waiting, the
agent cannot see them all, while MLFQ's aging rule still applies to
the full queue regardless of size.

We address this by retraining from scratch on Poisson arrival streams
($\lambda\sim U(0.01,0.08)$, 300 completions per episode). Retraining
under the new arrival distribution is necessary---the policy learned
on fixed five-task episodes does not transfer directly. When more
than five tasks are waiting, we show the agent the five that have
waited longest, prioritizing visibility of the most delayed tasks.
Training required four stability fixes: TD target clipping $[-50,50]$
(prevents Q-value divergence), Huber loss, frequent target network
sync (every 5 episodes), and a reduced learning rate ($10^{-4}$).

\begin{table}[h]
\centering
\caption{\omegavarn\ vs.\ MLFQ on variable-depth Poisson arrivals.}
\label{tab:w15}
\small
\begin{tabular}{lrr}
\toprule
Agent & MCT (mean$\pm$std, s) & Starv.\,(\%) \\
\midrule
MLFQ baseline & $22.06\pm13.36$ & 37.5 \\
\omegavarn\ seed 42 & $21.97\pm13.81$ & 40.0 \\
\omegavarn\ seed 123 & $20.57\pm11.96$ & 47.5 \\
\omegavarn\ seed 456 & $\mathbf{19.43\pm11.83}$ & 42.0 \\
\omegavarn\ mean$\pm$std & $20.66\pm1.03$ & 43.2 \\
\bottomrule
\end{tabular}
\end{table}

\omegavarn\ achieves lower mean MCT than MLFQ across all three seeds
(mean improvement: 1.40\,s; best seed: $-2.63$\,s). The effect is
modest relative to the high workload variance in this setting
(episode std $\approx12$--14\,s), so these gains should be
interpreted carefully. Starvation increases compared to MLFQ (43.2\%
vs.\ 37.5\%), consistent with the earlier observation that reward
shaping alone cannot enforce hard scheduling guarantees
(Section~\ref{sec:pbrs}).

The core bottleneck in this setting is the fixed-size input window
rather than the architecture itself. Scaling to a full transformer
architecture~\cite{vaswani2017} would allow the policy to handle
variable-length queues more naturally, without the truncation
heuristic, and is the most direct path to addressing this limitation.

%% ============================================================
\section{Discussion}
%% ============================================================

\subsection{Why Value-Delta Works}

We originally designed the value-delta reward as an urgency-aware
signal---prioritize tasks losing value fastest. It turned out to also
act as an implicit throughput signal. Tasks early in their decay
curves incur large value loss per unit of delay, so the agent learns
to schedule them quickly. This induces behavior qualitatively similar
to SRPT: reducing delay for high-loss tasks also reduces their
overall completion time. The two signals are not equivalent---SRPT
depends on remaining processing time while value-delta depends on
decay rate and current delay---but they correlate closely enough on
realistic workloads that the agent approximates SRPT ordering without
being told to.

More precisely, the agent learns to minimize aggregate value loss
across the whole queue at each step. A task that is currently
decaying fast contributes a large loss to that aggregate whether or
not it gets scheduled---so the agent learns to clear high-loss tasks
promptly, which is functionally similar to finishing short jobs first.

When we decompose this reward into separate urgency and completion
components (the W11 series, Section~\ref{sec:2x2}), both objectives
degrade. This suggests that the coupling between urgency and
completion efficiency is what makes the signal work. Separating them
into explicit additive terms breaks the implicit coordination that
the unified signal provides---each component can then be optimized
against the other rather than alongside it.

\subsection{The Starvation Gap}

\omegadqn\ at $\omega_s=0.7$ achieves 32.5\% starvation compared to
MLFQ's 36.0\%---a 3.5 percentage point reduction. This margin is
modest, and its magnitude depends on the specific threshold used in
our starvation metric (Section~\ref{sec:vrfi}). These are also
single-seed results for \omegadqn; we did not run a full multi-seed
sweep at every $\omega_s$ value, so the margin may shift with
different random seeds.

The underlying reason the gap is small connects to the earlier
finding in Section~\ref{sec:pbrs}: MLFQ enforces starvation
prevention through a hard aging rule that promotes any task waiting
beyond a fixed threshold regardless of other queue conditions. Our
approach uses reward shaping, which creates a soft incentive the
agent can trade off against throughput. In our experiments, that
tradeoff never fully closed the gap to zero, even at $\omega_s=0.7$
where the starvation signal is strong. The reduction is real, but it
does not carry the same guarantee as a hard rule.

\subsection{Limitations}

\begin{enumerate}[itemsep=2pt]

\item \textbf{Scale.} We evaluate on small instances: five processes,
  300-task episodes, single core. Real schedulers operate at orders
  of magnitude larger scale with far more concurrent processes and
  workload complexity.

\item \textbf{Fixed input size.} All models operate on a fixed window
  of five visible tasks, requiring truncation when queues grow
  longer. This limits performance under higher load---as the
  variable-N experiments show---and is the primary motivation for
  variable-length architectures as future work.

\item \textbf{Oracle burst.} \vadqn\ uses true remaining burst time,
  which is unavailable in real kernels. Removing it increases MCT by
  approximately 4.4\,s relative to the best checkpoint; the effect
  on the replicated mean is smaller (Section~\ref{sec:2x2}).

\item \textbf{Value model dependence.} Scheduling behavior depends on
  the choice of value curve parameters ($\tau$, floor, base). These
  encode assumptions about task urgency that are not learned from
  data. Different specifications will produce different policies, and
  there is no ground truth for what the right curves should be.

\item \textbf{Curve assignment.} Tasks are labeled steep or smooth by
  class. In a multi-tenant system, those labels are a policy choice
  with direct fairness consequences---whoever assigns them controls
  which tasks are treated as urgent. This is not a technical problem
  but a governance one.

\item \textbf{Trace truncation.} We restrict evaluation to tasks
  below the 75th percentile duration to limit burst ratio skew
  during training. This likely underrepresents long-tail workload
  behavior and may favor our method on the easier portion of the
  distribution.

\item \textbf{Seed count.} Three seeds confirms qualitative trends
  but limits statistical confidence, particularly for the single-seed
  \omegadqn\ frontier results and the \omegavarn\ starvation numbers.

\end{enumerate}

\subsection{Who Sets the Curves?}

Steep versus smooth assignment follows task class in our experiments.
In a shared system, the choice of $\tau$ directly determines how
costly delay is for each task. A smaller $\tau$ produces faster
early value decay, giving that task higher scheduling priority under
our framework. A larger $\tau$ produces slower decay and lower
effective priority.

This makes curve assignment a fairness question. Different $\tau$
values induce systematically different scheduling outcomes across
tenants---a tenant assigned a low $\tau$ receives better service not
because their tasks are inherently more urgent, but because the
parameter says they are. In a multi-tenant system, whoever controls
these parameters implicitly controls urgency, and by extension,
resource allocation.

Deploying value-aware scheduling therefore requires a governance
mechanism for assigning and auditing curve parameters---for example,
per-tenant policies, administrative quotas on $\tau$ ranges, or
third-party auditing of label assignments. We do not address this
here, but it is a prerequisite for any production deployment of this
approach.

%% ============================================================
\section{Conclusion}
%% ============================================================

We built and tested fifteen agent variants over twelve weeks on the
Alibaba 2018 cluster trace. Here is what we found.

\textbf{Local value-based scoring is insufficient.} The Quantum
Scoring Monotonicity Lemma shows that the two natural value-drop
objectives collapse to extreme quanta when applied per task in
isolation. Three ablations confirm this empirically. Attention
provides a practical way to address this by giving each task a view
of the whole queue before scoring.

\textbf{Existing fairness metrics miss urgency differences.} \vrfi\
captures value-rate inequality that JFI and SDV both report as fair.
The falsification case is concrete: identical wait times, identical
slowdowns, yet a $2\times$ difference in how fast tasks lose value.

\textbf{In our setting, reward structure is the primary lever.}
Replacing value-delta with a composite reward significantly degrades
performance. Removing oracle burst has a smaller effect. This
identifies reward design as the more important design choice between
the two.

\textbf{One model, adjustable at runtime.} \omegadqn\ improves both
throughput and starvation relative to MLFQ at $\omega_s=0.7$, with a
3.5 percentage point reduction in starvation and a modest MCT
improvement. The tradeoff is adjustable without retraining.

\textbf{Variable queues can be handled with retraining.} \omegavarn\
achieves lower mean MCT than MLFQ across seeds after training on
Poisson arrivals. The effect is modest relative to workload variance,
but it shows the approach is not limited to fixed-size settings.

\subsection*{What to Do Next}

Four directions worth pursuing: (1) train \omegadqn\ on Poisson
arrivals to combine preference conditioning with variable queue
depth; (2) learn burst prediction from observable features to reduce
the oracle gap; (3) use a full transformer~\cite{vaswani2017} to
handle variable-length queues natively, removing the five-slot
bottleneck; (4) explore whether large pretrained models can provide
useful scheduling heuristics without explicit RL training.

%% ============================================================
\bibliographystyle{plain}
\begin{thebibliography}{12}

\bibitem{codeRepo}
  V.~Upadhyay and S.~Brandt.
  \newblock Learning to Schedule: Code and Experiments.
  \newblock GitHub, 2025.
  \newblock \url{https://github.com/VedantUpadhyay/rl-cpu-scheduler}

\bibitem{mnih2015}
  V.~Mnih et~al.
  \newblock Human-level control through deep reinforcement learning.
  \newblock {\em Nature}, 518:529--533, 2015.

\bibitem{vaswani2017}
  A.~Vaswani et~al.
  \newblock Attention is all you need.
  \newblock {\em NeurIPS}, 2017.

\bibitem{film2018}
  E.~Perez, F.~Strub, H.~de~Vries, V.~Dumoulin, and A.~Courville.
  \newblock FiLM: Visual reasoning with a general conditioning layer.
  \newblock {\em AAAI}, 2018.

\bibitem{kernelOracle}
  S.~Y.~Kahu.
  \newblock KernelOracle: Predicting the Linux Scheduler's next move.
  \newblock {\em arXiv:2505.15213}, 2025.

\bibitem{alps}
  Y.~Fu et~al.
  \newblock ALPS: An adaptive learning, priority OS scheduler for
  serverless functions.
  \newblock {\em USENIX ATC}, 2024.

\bibitem{rlscheduler}
  D.~Zhang et~al.
  \newblock RLScheduler: An automated HPC batch job scheduler using RL.
  \newblock {\em SC}, 2020. doi:10.1109/SC41405.2020.00035.

\bibitem{doubleDQN}
  X.~Sun et~al.
  \newblock Dynamic OS scheduling using Double DQN.
  \newblock {\em arXiv:2503.23659}, 2025.

\bibitem{smartOS}
  S.~Goodarzy et~al.
  \newblock SmartOS: Towards automated learning and user-adaptive
  resource allocation in OS.
  \newblock {\em APSys}, 2021. doi:10.1145/3476886.3477519.

\bibitem{littlesLaw}
  J.~D.~C.~Little.
  \newblock A proof for the queuing formula: $L=\lambda W$.
  \newblock {\em Operations Research}, 9(3):383--387, 1961.

\bibitem{jainFairness}
  R.~Jain, D.~Chiu, and W.~Hawe.
  \newblock A quantitative measure of fairness and discrimination for
  resource allocation in shared computer systems.
  \newblock DEC TR-301, 1984.

\bibitem{alibaba}
  Alibaba Group.
  \newblock Alibaba cluster trace program, 2018.
  \newblock \url{https://github.com/alibaba/clusterdata}.

\bibitem{eevdf}
  T.~Natvig.
  \newblock EEVDF scheduler.
  \newblock {\em Linux Kernel Mailing List}, 2023.
  \newblock \url{https://lwn.net/Articles/925371/}.

\end{thebibliography}

\end{document}